\begin{titlepage}
\thispagestyle{empty}
\doublespacing
\textsc{\LARGE FOLIO - Il webmagazine di FedoraOnLine}\\[1.0cm]
\textsc{\Large rivista non professionale tematica e libera creata dai lettori di FedoraOnLine, scaricabile e dai contenuti forniti dagli utenti Fedora in italia e nel mondo.}\\[0.5cm]
\begin{figure}[htbp]
\centering
\includegraphics[scale=1.0]{immagini/logo/cc.png}
\end{figure}\\
{\bfseries Folio ed i suoi contenuti sono distribuiti con licenza creative commons, reperibile a questo link: http://creativecommons.org/licenses/by-sa/3.0/it/ }\\[0.4cm]
{\doublespacing\\}
\begin{tiny}
\begin{spacing}{1.0}
Questa rivista segue le linee guida dettate dal Fedora Project reperibile all'indirizzo  http://fedoraproject.org/wiki/Legal:Trademark\_guidelines:\\
{\bfseries Publications}\\
``It is permissible to use the Fedora trademarks in the title and content of a publication, provided that:
    the use is clearly in reference to the Fedora Project or its software
    the use does not imply sponsorship or endorsement by Red Hat or the Fedora Project
    Proper trademark symbols are used in connection with the Fedora Trademarks and the trademark attribution statement must appear as explained in Proper Trademark Use.'' \\
Fedora\textregistered and the Infinity design logo are trademarks of Red Hat, Inc.
\end{spacing}
\end{tiny}
\rule{\textwidth}{0.5mm}\\
{\doublespacing\\}
\vfill
\noindent
\begin{minipage}{0.4\textwidth}
\begin{flushleft} 
\singlespacing
\large
{\bfseries\color[cmyk]{0.5, 0, 1, 0}{\emph {Redazione:}}}\\
\bigskip
Gabriele \textsc{Trombini}\\
{\tiny mailga@fedoraonline.it} \\
\bigskip
Robert \textsc{Mayr}\\
{\tiny robyduck@fedoraonline.it} \\
\bigskip
Giuseppe \textsc{Delvecchio}\\
{\tiny virus@fedoraonline.it} \\
\bigskip
Mario \textsc{Santagiuliana}\\
{\tiny marios@fedoraonline.it} \\
\bigskip
Giuseppe \textsc{Raveduto}\\
{\tiny wardialer@fedoraonline.it} \\
\bigskip
{\bfseries\color[cmyk]{0.5, 0, 1, 0}{\emph {Grafica:}}}\\
\bigskip
Cristian \textsc{Pozzessere}\\
{\tiny cristianpozzessere@gmail.com} \\
\end{flushleft}
\end{minipage}
\begin{minipage}{0.6\textwidth}
\begin{flushright} 
\singlespacing
\large
{\bfseries\color[cmyk]{0.5, 0, 1, 0}{\emph {Collaboratori:}}}\\
\bigskip
Andrea \textsc{Veri}\\
{\tiny averi@fedoraproject.org}\\
\bigskip
{\bfseries\color[cmyk]{0.5, 0, 1, 0}{\emph {Interventi:}}}\\
\bigskip
Gianluca \textsc{Sforna}\\
{\tiny giallu@fedoraproject.org}\\ 
\end{flushright}
\end{minipage}
\newpage
\flushcolumns
\twocolumn
\singlespacing
\title{\bfseries\color[cmyk]{0.5, 0, 1, 0} IN QUESTO NUMERO}
\date{}
\maketitle
\thispagestyle{empty} 
\normalsize
\flushbottom
\section*{\color[cmyk]{1, 0.57, 0, 0.38}Editoriale {\scriptsize(pag.1)}} \label{sec:Prima}
\small{Spiegati i motivi di questo Webmagazine, fortemente voluto.}\\ {\tiny di Gabriele Trombini}
\section*{\color[cmyk]{1, 0.57, 0, 0.38}Il Fedora Project {\scriptsize(pag.3)}} \label{sec:Seconda}
Cos'è il Fedora Project? Ecco cosa, dove, come e perchè!\\ {\tiny di Gabriele Trombini}
%\noindent\ref{sec:Terza}\pageref{sec:Terza}
\section*{\color[cmyk]{1, 0.57, 0, 0.38}Intervista a: Gianluca Sforna {\scriptsize(pag.7)}} \label{sec:Terza}
Uno degli Ambassador tra i più attivi risponde alle nostre domande.\\ {\tiny di Gabriele Trombini}
%\noindent\ref{sec:Quarta}\pageref{sec:Quarta}
\section*{\color[cmyk]{1, 0.57, 0, 0.38}Fedora Inside {\scriptsize(pag.10)}} \label{sec:Quarta}
Presentazione di Andrea Veri; la nostra guida all'interno del progetto Fedora.\\  {\tiny di Gabriele Trombini}
%\noindent\ref{sec:Quinta}\pageref{sec:Quinta}
\section*{\color[cmyk]{1, 0.57, 0, 0.38}FedoraOnLine {\scriptsize(pag.11)}} \label{sec:Quinta}
Cronistoria e retroscena del passaggio da FOL 1 a FOL 2.\\ {\tiny di Gabriele Trombini}
%\noindent\ref{sec:Sesta}\pageref{sec:Sesta}
\section*{\color[cmyk]{1, 0.57, 0, 0.38}Primi passi {\scriptsize(pag.14)}} \label{sec:Sesta}
Linea di comando: impariamo a conoscerla (prima parte).\\ {\tiny di Gabriele Trombini}
%\noindent\ref{sec:Settima}\pageref{sec:Settima}
\section*{\color[cmyk]{1, 0.57, 0, 0.38}Sistema base {\scriptsize(pag.19)}} \label{sec:Settima}
Panoramica sulla gestione dei servizi (prima parte).\\ {\tiny di Gabriele Trombini}
%\noindent\ref{sec:Decima}\pageref{sec:Decima}
\section*{\color[cmyk]{1, 0.57, 0, 0.38}Sistema base {\scriptsize(pag.23)}} \label{sec:Ottava}
Dracut, cosa è ed a cosa serve.\\ {\tiny di Gabriele Trombini}
%\noindent\ref{sec:Ottava}\pageref{sec:Ottava}
\section*{\color[cmyk]{1, 0.57, 0, 0.38}Sistema avanzato {\scriptsize(pag.26)}} \label{sec:Nona}
La creazione di un rpm, struttura di base (prima parte).\\ {\tiny di Gabriele Trombini}
%\noindent\ref{sec:Undicesima}\pageref{sec:Undicesima}
\section*{\color[cmyk]{1, 0.57, 0, 0.38}Sistema avanzato {\scriptsize(pag.29)}} \label{sec:Decima}
Introduzione al Kernel Fedora (prima parte).\\ {\tiny di Giuseppe Delvecchio}
%\noindent\ref{sec:Nona}\pageref{sec:Nona}
\section*{\color[cmyk]{1, 0.57, 0, 0.38}Sistema e sicurezza {\scriptsize(pag.32)}} \label{sec:Undicesima}
Cosa è e come funziona il keyserver Fedora, generazione di chiavi gpg.\\ {\tiny di Andrea Veri}
%\noindent\ref{sec:Dodicesima}\pageref{sec:Dodicesima}
\section*{\color[cmyk]{1, 0.57, 0, 0.38}Notizie {\scriptsize(pag.35)}} \label{sec:Dodicesima}
Sveliamo le features della prossima release: Fedora 17 (Beefy Miracle).\\ {\tiny di Robert Mayr}
\end{titlepage}

